% Generated draft: requires mathematical and source review.
\documentclass[11pt]{article}
\usepackage[margin=25mm]{geometry}
\usepackage{fontspec}
\setmainfont{Latin Modern Roman}
\usepackage{amsmath,amssymb,mathtools,unicode-math}
\setmathfont{Latin Modern Math}
\usepackage{enumitem,needspace,xurl,hyperref}
\setlength{\parindent}{0pt}
\setlength{\parskip}{5pt}
\setlength{\emergencystretch}{3em}
\newcommand{\N}{\mathbb{N}}
\newcommand{\Z}{\mathbb{Z}}
\newcommand{\Q}{\mathbb{Q}}
\newcommand{\R}{\mathbb{R}}
\newcommand{\C}{\mathbb{C}}
\newcommand{\F}{\mathbb{F}}
\newcommand{\A}{\mathbb{A}}
\newcommand{\PP}{\mathbb{P}}
\newcommand{\E}{\mathbb{E}}
\newcommand{\eps}{\varepsilon}
\newcommand{\dd}{\,\mathrm{d}}
\begin{document}
\section*{UnsolvedMath ID: 30004540}
\textbf{Source problem code:} catalogue entry
\subsection{Definitions and mathematical statement}
\textbf{Unrecoverable formulation.} The supplied question is: ``What is the asymptotic behavior of the persistence probability for an infinitely correlated Gaussian stationary process, and how does it depend on whether its correlation function is summable?'' The record does not define ``infinitely correlated.'' Its background also identifies a malformed phrase, ``Pa-infinity-correlated.'' Neither phrase specifies a mathematical hypothesis.

For clarity, the following describes possible notation only; it does not recover the intended problem. If the intended process is real-valued, write \(X=(X_t)_{t\in I}\), where the index set \(I\) must be specified. A Gaussian process has jointly Gaussian finite-dimensional distributions; stationarity means that these distributions are invariant under every common time translation preserving the index set. Its mean is constant, and its covariance function is
\[
 c(h)=\E[(X_{t+h}-\E X_{t+h})(X_t-\E X_t)].
\]
If \(c(0)>0\), its normalized correlation function is \(\rho(h)=c(h)/c(0)\). Centering and nondegeneracy are not supplied assumptions.

For a specified level \(u\), a possible persistence probability is
\[
 p_{X,u}(H)=\mathbb P\bigl(X_t\geq u\text{ for every }t\in I\cap[0,H]\bigr),
\]
provided this event is measurable. This is only an illustrative convention: the source does not specify the level, the inequality, the time domain, or whether persistence concerns \(X\) itself. In continuous time, suitable measurability or path assumptions would also be needed.

In discrete time, absolute summability might mean \(\sum_{k\in\Z}|\rho(k)|<\infty\); convergence of a signed correlation series is a different condition. In continuous time, an integrability condition such as \(\int_{\R}|\rho(t)|\,\dd t<\infty\) would be a different possible interpretation. The record does not select among these conditions or specify a correlation-sign hypothesis.

\textbf{Intended target and unresolved qualification.} The extraction asks for persistence asymptotics as the time horizon tends to infinity and a dependence on correlation summability. However, the supplied statement-verification field says that the report passage introduces results on weighted sums of correlated Gaussian stationary processes and subsequently states how the answer depends on summability. The weighted process, weights, normalization, hypotheses, and intended asymptotic quantity are absent. Thus no precise asymptotic assertion or separate open problem can be stated from this record without additional source information. The notation above must not be treated as a replacement formulation.
\subsection{Short English statement}
The record concerns persistence asymptotics associated with correlated stationary Gaussian processes, apparently through weighted sums. The process, weights, persistence event, and summability condition are unspecified. The supplied verification describes a lead-in to stated results, while the registry retains an \texttt{open} label; a precise open problem cannot be recovered.
\subsection{Sources}
\begin{itemize}
\item Local UnsolvedMath registry: supplied complete source record, permanent ID \(30004540\), source problem code \texttt{catalogue entry}, in the Oberwolfach Reports Open Problems collection. The record supplies proposed year \(2020\), but no named proposer.
\item Supplied report citation: \emph{Stochastic Processes under Constraints} (2021), Oberwolfach Reports, DOI \texttt{10.4171/owr/2020/32}, pp.~1601--1656. \url{https://doi.org/10.4171/owr/2020/32}. This is the report attributed by the record; no problem-specific page or passage is supplied.
\item Supplied literature source labeled ``Persistence spectral methods'': \url{https://arxiv.org/abs/1709.00204}. The record supplies no authors, full bibliographic title, or theorem identifying its relevance to a recovered formulation.
\item Literature assessment dated exactly \texttt{2026-08-22}: supplied dated curation, not independently verified. It labels the item open and says that the malformed extracted phrase does not identify a precise model or theorem, and that existing persistence results do not settle the unspecified generalization. The background expressly requests independent verification of time-sensitive status and cited claims.
\item Supplied statement-verification qualification: the report passage is described as a lead-in to results about weighted sums, followed by an answer depending on correlation summability. The record marks the statement \texttt{unrecoverable}, supplies no clean statement, and retains the original extraction for provenance. Its title is ``Resolved Weighted-Gaussian-Persistence Lead-In,'' while its status remains \texttt{open}. These are supplied registry descriptions, not independently verified mathematical conclusions.
\end{itemize}
\end{document}
